\documentclass[10pt]{article}
\usepackage[a4paper,margin=26mm]{geometry}
\usepackage{amsmath,amssymb,amsthm,mathrsfs,hyperref}
\newtheorem{theorem}{Theorem}
\newtheorem{lemma}{Lemma}
\newcommand{\vol}{\operatorname{vol}}
\newcommand{\MA}{\operatorname{MA}}
\title{Nef model moments and a mixed complete-moving-system limit}
\author{An unreviewed algebraic bridge for the big-class equation}
\date{6 October 2026}
\begin{document}
\maketitle
\begin{abstract}
We reconstruct a uniform raw moment estimate on actual nef models,
without an ample lower bound or probability normalization. Simultaneous
complete systems recover the mixed ordinary positive products of finitely
many big compactifications. For a shared generic bundle their actual
horizontal factors agree, so a difference of two mixed positive products
is the limit of an auxiliary vertical moment. Conditional on
the separately stated genuine energy foundations, that limit has a
uniform finite-energy bound. Identification with integration of the
original relative potential against the original output measure is a
separate unresolved step. No unbounded operator or inverse equation is
claimed by this paper.
\end{abstract}

\paragraph{Domain and boundary.}
All varieties are smooth integral projective over an algebraically
closed trivially valued field of characteristic zero. All bundles,
models, Cartier divisors and section systems are actual; rational
coefficients are cleared in their degrees. Nef model potentials are
bounded functions on the full original analytic carrier, given by
vertical Cartier divisors on actual equivariant projective models.
The finite model intersection calculus, its raw positive mixed measures,
and its vertical Hodge inequality are ordinary algebraic background
\cite{BJ}. All central multiplicities are retained. Finite model
statements below require neither a nonample decreasing-closure theorem
nor a target-dependent inverse operator.

\section{A raw estimate on a nef generic bundle}
Let $A$ be an actual nef rational bundle on $Z$, $\dim Z=n\ge1$, and
write $V_A=A^n\ge0$. For a bounded relatively nef model potential
$\psi$ write $g_\psi=\psi(v_0)=\sup\psi$. Define its raw relative
intersection energy $e_A$ with trivial reference, and
\[
 J_A(\psi)=g_\psi V_A-e_A(\psi)\ge0,
 \qquad \mu_A(\psi)=\MA_{A,\mathrm{raw}}(\psi).
 \tag{1}
\]
Thus $\mu_A(\psi)(Z^{\rm an})=V_A$ and
$\mu_A(0)=V_A\delta_{v_0}$. In particular we never divide by $V_A$.
Generic value is the supremum by the finite section/model convention;
it also follows by passing through an arbitrarily small actual ample
zero-metric addition. The model measure is invariant under constants.

\begin{theorem}\label{moment}
There exists $C_n$ depending only on dimension such that for all
bounded relatively nef model potentials $h,\psi$ on $A$,
\[
 \begin{split}
 \int(g_h-h)\,d\mu_A(\psi)
 &\le C_n J_A(h)^{\alpha_n}J_A(\psi)^{1/2}
       \max\{J_A(h),J_A(\psi)\}^{1/2-\alpha_n},\\
 &\hspace{35mm}\alpha_n=2^{-n}.
 \end{split}                                           \tag{2}
\]
A zero right side is interpreted by continuity, including $V_A=0$.
If $h\le0$ then
\[
 \int(-h)\,d\mu_A(\psi)
 \le(-g_h)V_A+C_n J_A(h)^{\alpha_n}J_A(\psi)^{1/2}
       \max\{J_A(h),J_A(\psi)\}^{1/2-\alpha_n}.          \tag{3}
\]
The same constant works for all varieties, actual bundles, models and
central multiplicities in this signature.
\end{theorem}

\paragraph{The finite Hodge mechanism.}
We give the algebraic derivation, to expose the uniformity needed in a
moving generic class. Write $T_\xi=(A,\xi)$ as notation for model
intersection slots, not for an unbounded current. On a common model,
for nef slots $T_1,\ldots,T_{n-1}$ the vertical Hodge inequality gives
\[
 \|f\|_{T_1\cdots T_{n-1}}^2
 :=-(0,f)^2T_1\cdots T_{n-1}\ge0.                     \tag{4}
\]
Polarizing (4) proves Cauchy--Schwarz. Define
\[
 d_A(\xi,\eta)=\max_{0\le j<n}
            \|\xi-\eta\|_{T_\xi^jT_\eta^{n-1-j}}^2.
\]
The finite telescoping formula and successive nef replacements give
\[
 I_A(\xi,\eta)=\int(\xi-\eta)(\mu_A(\eta)-\mu_A(\xi))
     =\sum_{j=0}^{n-1}\|\xi-\eta\|_{T_\xi^jT_\eta^{n-1-j}}^2.
 \tag{5}
\]
For $\eta=0$, the same telescoping calculation yields
\[
 J_A(\xi)=\frac1{n+1}\sum_{j=0}^{n-1}(j+1)
                         \|\xi\|_{T_\xi^jT_0^{n-1-j}}^2.
 \tag{6}
\]
Hence $d_A(\xi,0)$, $I_A(\xi,0)$ and $J_A(\xi)$ are comparable
with constants depending only on $n$. These formulas are raw and
hold even when $V_A=0$.

For completeness, the finite change-of-slots argument is as follows.
Set $f=\xi-\xi'$, $\tau=(\xi+\xi')/2$, $a=d_A(\xi,\xi')$,
$b=\max\{d_A(\xi,\eta),d_A(\xi',\eta)\}$ and
\[
 b_j=\|f\|_{T_\eta^jT_\tau^{n-1-j}}^2.
\]
Expansion gives $b_0\le C_na$. Splitting $f$ into its two differences
with $\eta$ gives $b_{n-1}\le4b$. For $j<n-1$, splitting the difference
$b_{j+1}-b_j$ into the two slots $T_\xi,T_{\xi'}$, Cauchy--Schwarz
and $T_\xi,T_{\xi'}\le2T_\tau$ give
\[
 b_{j+1}-b_j\le C_n\sqrt{b b_j}.                     \tag{7}
\]
Here $d_A(\tau,\eta)\le C_nb$ follows from the finite energy
concavity and (6)'s analogous two-reference identity. If $a\le b$,
induction in (7) gives
\[
 b_j\le C_n a^{2^{-j}}b^{1-2^{-j}}.
\]
If $a>b$, use $b_{n-1}\le4b$ instead. Thus the final-slot estimate has
exponent $2^{1-n}$. For $n=1$ the product of zero slots is interpreted
literally and the assertion follows from the seminorm triangle inequality.

Using $\tau=(\xi+\eta)/2$ and splitting each difference through a
third potential, the same estimate gives the quasi-triangle inequality
\[
 d_A(\xi,\eta)\le C_n\bigl(d_A(\xi,\zeta)+d_A(\eta,\zeta)\bigr).
\]
One can see the absorption explicitly: if the left side is the larger
quantity, the change-of-slots bound has the form
$d\le C_n b^{2^{1-n}}d^{1-2^{1-n}}$, hence $d\le C_n'b$.
A mixed product of $n-1$ nef slots is bounded by a dimension constant
times the $(n-1)$st power of their average. Apply the final-slot
estimate and the quasi-triangle inequality to this average to obtain
\[
 \|\xi-\xi'\|_{T_{\eta_1}\cdots T_{\eta_{n-1}}}^2
 \le C_n d_A(\xi,\xi')^{2^{1-n}}
       \max\{d_A(\xi,0),d_A(\xi',0),d_A(\eta_i,0)\}^{1-2^{1-n}}.
 \tag{8}
\]
These are the finite nef identities of \cite[Section 3.3]{BJ}; no
ample energy extension is being applied to a merely nef class.

\paragraph{Proof of Theorem~\ref{moment}.}
Shift $h,\psi$ to generic value zero. Then $\int h\,d\mu_A(0)=0$.
Expanding the difference of the $n$th powers gives
\[
 \int h\,d(\mu_A(\psi)-\mu_A(0))
   =\sum_{j=0}^{n-1}(0,h)(0,\psi)T_\psi^jT_0^{n-1-j}. \tag{9}
\]
Cauchy--Schwarz bounds the $j$th term by
$\|h\|_{T_\psi^jT_0^{n-1-j}}\|\psi\|_{T_\psi^jT_0^{n-1-j}}$.
The second factor is at most $d_A(\psi,0)^{1/2}$.
By (8) the first factor is at most
\[
 C_n d_A(h,0)^{2^{-n}}
             \max\{d_A(h,0),d_A(\psi,0)\}^{1/2-2^{-n}}.
\]
Use (6) and sum the finite terms. This proves (2); undoing the constant
shift proves (3). Continuity of finite intersection polynomials handles
all zero seminorms and real model coefficients. \qed

\section{Simultaneous complete systems recover mixed positive products}
Let $Y$ be smooth projective of dimension $d$ and let
$N_1,\ldots,N_p$ be actual big rational bundles, $1\le p\le d$.
On a simultaneous resolution of the full $|mN_i|$ systems write
\[
 \mu^*N_i=C_{m,i}+G_{m,i}/m,
       \qquad mC_{m,i}\text{ globally generated},\quad G_{m,i}\ge0.
\]
Take cofinal factorial degrees. Each $C_{m,i}$ increases in the
pseudoeffective order on common higher models. Then
\[
 C_{m,1}\cdots C_{m,p}
                  \longrightarrow\langle N_1\cdots N_p\rangle
 \tag{10}
\]
in the weak Weil numerical-class topology.

Here is a direct cofinality proof. Given finitely many nef competitors
$\beta_i\le\mu^*c_1(N_i)$ on a common model, choose a sufficiently
small actual ample $Q_0$ so all $\mu^*N_i-Q_0$ are big. Perturb
$\beta_i$ to $Q_i=(1-\epsilon)\beta_i+\epsilon c_1(Q_0)$.
For rational defining competitors these are actual ample rational
bundle classes; choose actual representatives. Their actual differences
from $\mu^*N_i$ are big. For one sufficiently divisible $q$ each such
difference has an effective divisor $R_i$ and $qQ_i$ is free:
\[
 q\mu^*N_i=qQ_i+R_i,\qquad R_i\ge0.
\]
The generated subseries implies $G_{q,i}\le R_i$, so $C_{q,i}\ge Q_i$.
Nef-product monotonicity gives $\prod C_{q,i}\ge\prod Q_i$.
Let $\epsilon\downarrow0$. Directed compactness and the defining least
upper bound prove (10) \cite{BFJ}. Real competitors are reached by the
usual ample approximation and closedness; actual section identities
are used only at rational stages. Products with repeated factors use
the same complete system and need no independent choices. Degree-$d$
products are numbers, so (10) is ordinary scalar convergence there.

\section{A formal vertical moment on the original model data}
Let $X$ have dimension $n$, $L$ be actual big, $V=\vol(L)>0$.
Fix rational model obstacles $f_D,f_H$ on one smooth equivariant
compactification $Y\to X\times\mathbf P^1$, with vertical Cartier
$D,H$ and $F=\sum b_EE$. A single integer $c$ exceeds the coefficient
bounds for both obstacles and zero. For $Z\in\{0,D,H\}$ set
\[
 \alpha_Z=\pi_X^*L+Z+cF.
\]
These are actual big bundles. Indeed their complete relative lattices
have weights in $[-mM,mM]$ for a common $M<c$, and their full absolute
bases are $t^{-(w_{mj}+mc)+r}s_{mj}$ for $0\le r\le w_{mj}+mc$.
Thus $h^0(Y,m\alpha_Z)\ge(m(c-M)+1)h^0(X,mL)$.
Their restrictions give the full relative lattice and the full generic
system at infinity.

Resolve all three absolute complete systems at the same degree $m$,
and also the complete original $|mL|$. Write its generic moving
factorization $p_m^*L=A_m+T_m$, with $mA_m$ free and $T_m\ge0$.
Every absolute section has the same horizontal factor $mT_m$.
After its removal each residual ideal is the unit ideal off time zero,
since the full generic moving system is free, including at infinity.
Consequently the three complete moving bundles have the actual form
\[
 C_{m,Z}=\rho^*A_m+Z+cF-G_{m,Z}^{\rm vert}/m,
       \qquad C_{m,0}=\rho^*A_m+cF.                    \tag{11}
\]
They are globally generated and nef. Removing $cF$ yields bounded
relatively nef model potentials $\psi_{m,Z}$ on the SAME actual $A_m$.
All identities are on this auxiliary carrier; no identification of its
supported points with the original supported points is used.

The finite vertical moment is exactly
\[
 M_m(H,D):=\int(-\psi_{m,H})\,d\mu_{A_m}(\psi_{m,D})
            =(C_{m,0}-C_{m,H})C_{m,D}^{n}.            \tag{12}
\]
This follows from the model component formula with every $b_E$:
the vertical divisor's coefficient at $v_E$ is divided by $b_E$,
while the measure coefficient is multiplied by $b_E$. Projection to
higher common models preserves the actual intersection identity.
By (10) the scalar limit exists and equals
\[
 M_{\rm alg}(H,D)
 :=\langle\alpha_0\alpha_D^n\rangle
                     -\langle\alpha_H\alpha_D^n\rangle.
 \tag{13}
\]
This is a difference of mixed ORDINARY positive products. It is not an
assertion that the positive product is multilinear in its big-class
arguments, nor a definition of an unbounded original current product.
Since (12) is independent of any sufficiently large common time twist,
its limit (13) has the same independence. This proves only the indicated
difference's independence; the two individual products vary with $c$.

\paragraph{Conditional finite-energy bound.}
Suppose the separate genuine closure and section-volume foundations
have been proved. In particular let $u_Z=P_\theta(f_Z)$ be normalized
with $u_Z(v_0)=0$, let $E(u_Z)$ be their genuine raw energy, and assume
the already proposed full-system identities
\[
 \phi_{m,Z}\circ p_m^{\rm an}=q_{T_m}+\psi_{m,Z},\quad
 g_{m,Z}:=\psi_{m,Z}(v_0)\to0,\quad
 e_{A_m}(\psi_{m,Z})\to E(u_Z),\quad A_m^n\to V.       \tag{14}
\]
For normalized model envelopes the finite potentials are nonpositive;
their generic values are their suprema. The raw model costs therefore
satisfy $J_{A_m}(\psi_{m,Z})\to B(u_Z):=-E(u_Z)$.
Applying (3) before taking the moving degree limit gives
\[
 0\le M_{\rm alg}(H,D)
 \le C_n B(u_H)^{\alpha_n}B(u_D)^{1/2}
           \max\{B(u_H),B(u_D)\}^{1/2-\alpha_n}.         \tag{15}
\]
The generic error $(-g_{m,H})A_m^n$ tends to zero. Thus there is a
uniform algebraic moment bound under exactly the foundations (14).
We have kept the changing generic volume and the generic-value error;
there is no ample lower-volume or countability restriction.

\section{The remaining identification is substantive}
For the original big-class output measure $\mu(u_D)$, the equation
programme still needs to prove
\[
 M_{\rm alg}(H,D)
       =\int_{\{V_\theta> -\infty\}}(V_\theta-u_H)\,d\mu(u_D)
 \tag{16}
\]
and that the omitted reference-pole set has zero output mass.
Equation (16) is \emph{not proved here}. The left side is an auxiliary
moving-series limit; the right side is integration of an actual
original-carrier relative difference. Weak convergence of measures and
pointwise convergence of changing bounded test functions alone cannot
justify (16). Even the shared horizontal-factor cancellation in (11)
does not supply uniform original continuous-test errors or weighted
convergence near poles. An original-carrier negligible-discrepancy or
weighted-identification proof must do that work.

After (16), extension to every unbounded genuine finite-energy test,
strong fixed-target integration, strict comparison and actual arbitrary
target attainment are further required. The finite estimate (2) has
been algebraically reconstructed, and (10)--(13) give its precise moving
limit. They do not settle any of those subsequent steps. This unreviewed
supporting draft closes no original problem and changes no accepted
paper count.

\begin{thebibliography}{9}
\bibitem{BJ} S. Boucksom and M. Jonsson,
\emph{Global pluripotential theory over a trivially valued field},
\href{https://arxiv.org/pdf/1801.08229v4}{arXiv:1801.08229v4},
Section 3.3, especially Lemmas 3.24, 3.32--3.33 and Corollary 3.34:
finite model nef calculus. Lemma 7.30's ample finite-energy extension
is compared only after reconstructing the finite argument above.
Proposition 9.19's strong target integration requires its prior
quantitative weighted estimates; it cannot simply be imported into a
big class. Hodge-index foundations remain ordinary primary background.
\bibitem{BFJ} S. Boucksom, C. Favre and M. Jonsson,
\emph{Differentiability of volumes of divisors and a problem of Teissier},
\href{https://arxiv.org/html/math/0608260v2}{arXiv:math/0608260v2},
Lemma 2.2 and Proposition 2.3 with proofs, Definition 2.5,
Lemmas 2.6--2.7 with proofs and Proposition 2.13.
The characteristic-zero algebraic domain and genuine birational
positive-product topology are retained.
\end{thebibliography}
\end{document}
